\chapter*{Conclusion générale}
\addcontentsline{toc}{chapter}{Conclusion générale}

Ce projet de fin d’études réalisé dans le cadre de notre stage au sein de l’entreprise Virtual Dev nous a permis de mettre en pratique les connaissances acquises durant notre parcours universitaire et de développer davantage nos compétences techniques et professionnelles dans le domaine du développement logiciel.

L’objectif principal de ce projet était la conception et le développement d’une plateforme intelligente de gestion et de validation des factures fournisseurs basée sur les technologies web, l’OCR et l’intelligence artificielle. Cette solution a été proposée afin d’automatiser le traitement des factures, réduire les tâches manuelles et améliorer l’efficacité du processus de validation financière.

Tout au long de ce travail, nous avons commencé par étudier l’existant et identifier les différentes limites du traitement manuel des factures, notamment la perte de temps, les erreurs humaines et le manque de traçabilité. Ensuite, nous avons analysé les besoins fonctionnels et non fonctionnels du système afin de proposer une architecture adaptée et performante.

La réalisation de ce projet nous a également permis d’utiliser plusieurs technologies modernes telles que React pour le frontend, FastAPI pour le backend, PostgreSQL pour la gestion de la base de données ainsi que les outils OCR et IA pour l’extraction et l’analyse automatique des données. Nous avons aussi adopté la méthodologie agile SCRUM afin d’assurer une meilleure organisation du travail et un suivi efficace de l’avancement du projet.

La plateforme développée offre plusieurs fonctionnalités importantes, notamment l’importation des factures, l’extraction automatique des données, la validation financière, la gestion des utilisateurs, le suivi des actions et la visualisation des statistiques via des tableaux de bord interactifs. Cette solution permet ainsi d’améliorer la rapidité, la fiabilité et la sécurité du traitement des factures au sein de l’entreprise.

Ce projet représente une expérience très enrichissante aussi bien sur le plan technique que professionnel. Il nous a permis de travailler sur un projet concret, de résoudre plusieurs problèmes techniques et de découvrir de nouvelles technologies utilisées dans le domaine des systèmes intelligents.

Comme perspectives d’amélioration, il serait intéressant d’ajouter de nouvelles fonctionnalités telles qu’un système avancé de notification, une analyse intelligente plus poussée des anomalies, ainsi qu’une intégration avec d’autres systèmes de gestion d’entreprise afin d’améliorer davantage les performances et les capacités de la plateforme.

En conclusion, ce projet constitue une étape importante dans notre parcours académique et professionnel. Les résultats obtenus répondent aux objectifs fixés et représentent une base solide pour de futures améliorations et évolutions du système.